\documentclass[11pt]{article}
\usepackage[margin=1in]{geometry}
\usepackage{amsmath,amssymb,amsthm}

\newtheorem{definition}{Definition}
\newtheorem{proposition}{Proposition}

\title{Linear Algebra II\\\large Lecture 7: Eigenvalues and diagonalisation}
\author{Course notes}
\date{Winter term}

\begin{document}
\maketitle

\section{Eigenvalues}
Let $A \in \mathbb{R}^{n \times n}$. A scalar $\lambda$ is an eigenvalue of
$A$ when there is a non-zero vector $v$ with
\begin{equation}
  A v = \lambda v .
\end{equation}

\begin{definition}
The characteristic polynomial of $A$ is $p_A(\lambda) = \det(A - \lambda I)$.
\end{definition}

\begin{proposition}
$\lambda$ is an eigenvalue of $A$ if and only if $p_A(\lambda) = 0$.
\end{proposition}

\begin{proof}
$A v = \lambda v$ has a non-zero solution exactly when $A - \lambda I$ is
singular, that is, when its determinant vanishes.
\end{proof}

\section{Diagonalisation}
If $A$ has $n$ linearly independent eigenvectors $v_1, \dots, v_n$, collect
them as the columns of $P$. Then
\begin{equation}
  P^{-1} A P = \operatorname{diag}(\lambda_1, \dots, \lambda_n).
\end{equation}

\paragraph{Example.} For
$A = \begin{pmatrix} 2 & 1 \\ 1 & 2 \end{pmatrix}$ the eigenvalues are
$1$ and $3$, with eigenvectors $(1,-1)^\top$ and $(1,1)^\top$.

\section{Exercises}
\begin{enumerate}
  \item Show that similar matrices share their characteristic polynomial.
  \item Find a $2 \times 2$ matrix that is not diagonalisable over $\mathbb{R}$.
  \item Prove that a symmetric matrix has real eigenvalues.
\end{enumerate}

\end{document}
